\scanpage{205}
\mspage{\unsure{61}}
dans chacun des cas suivants

$1^\circ)$ $I$ est préordonné

$2^\circ)$ $\forall\ i$ de $I$, $\alpha_i^*\colon\mathbb{D}(I)\mathbin{\unsure{-|\to}}\mathbb{D}(I_{/i})$ est ess. surjectif.

% The transfer arrow in the left margin leads from the line above 2 degrees to the space below it.
\inserted{$3^\circ)$ $\forall\ i$ de $I$, \unsure{considérant|écrivant} $\alpha_i\colon I_{/i}\to I$, le foncteur $\alpha_{i!}\colon$ \inserted{$\mathbb{D}(I_{/i})\mathbin{-}\mathbb{D}(I)$} est conservatif}

$4^\circ)$ $\mathbb{D}=\mathrm{HOT}_{W_0}\approx\mathbb{D}_{\mathrm{Ens}}$.

Dém. L'équiv. de $a,a',a''$ est claire, soit $A$. De $=$ l'équivalence de $b,b',b''$, soit $B$.

Prouvons

\[
 B\Rightarrow A
\]

en supposant p. ex. l'existence de $u_*$. On admet tacitement, avec l'existence de $u_*\colon\mathbb{D}(I')\mathbin{-}\mathbb{D}(I)$, la morale du calcul des \deleted{foncteurs} fibres de $u_*(-)$. Donc $A$, via $a''$, équivaut à

\[
 \unsure{\mathbb{H}|H}^{\bullet}_{\mathbb{D}}(I'_{/i},u^*F)
 \xleftarrow{\text{\uuline{iso}}}
 H^{\bullet}_{\mathbb{D}}(I_{/i},F)\simeq F(i)
 \tag*{(\unsure{x|\ast})}
\]

pour tt $F$ de $\mathbb{D}(I)$, \unsure{tout} $i$ de $I$.

\deleted{Désignant par $F_i$ \unsure{(}}

Soit $\alpha_i\colon I_{/i}\mathbin{-}I$, $F_i=\alpha_i^*(F)$, on doit donc exprimer que pour tt $F_i$ \unsure{provenant} d'un $F$, on a

\[
 H^\bullet_{\mathbb{D}}(I_{/i},F_i)\xrightarrow{\mathrm{iso}}
 H^\bullet_{\mathbb{D}}(I'_{/i},(u_{/i})^*(F_i)).
 \tag*{(\unsure{xx|\ast\ast})}
\]

Or \deleted{dire que} sous réserve d'existence de $(u_{/i})_!$, dire que $(\unsure{xx|x})$ \unsure{est|$\to$} iso pour \uline{tt} $\unsure{F_i|F'}$ de $\mathbb{D}(I_{/i})$, signifie que \deleted{\unsure{a}}.

\[
 u_{/i}\colon I'_{/i}\longrightarrow I_{/i}
\]
% archive sheet number: 205
